\chapter{The Research Environment}\label{ch:pl:the-research-environment}

In 2013 three economists replicating an influential paper on public debt and growth obtained the authors' working spreadsheet. A coding
error in it excluded five countries --- Australia, Austria, Belgium, Canada and Denmark --- from the analysis; with that and two other
problems corrected, the average growth of countries with public debt above 90\% of output was 2.2\%, not the $-0.1\%$ published (Herndon,
Ash and Pollin). The economics was not what failed. What failed was an environment in which a calculation could not be rerun from the top
by anyone but its author, and not even by its author once the state of the spreadsheet had drifted from what it showed. A research
platform's notebooks have the same property, more so. This chapter audits a notebook that lies in the same way, reruns it from the top,
promotes its calculation to a library, and reproduces the corrected number from a clean directory.

\section{Notebooks and their hidden state}

\begin{definition}[Computational notebook]\label{def:pl:the-research-environment:notebook}
A \emph{computational notebook}\index{computational notebook} is a document of cells --- code, text and the outputs the code produced ---
run one cell at a time against a live interpreter (the kernel), in any order the user chooses; Jupyter's format stores it as JSON, with each
code cell's source, its saved outputs and the kernel's execution counter when it last ran (an integer, or null if it never ran).
\end{definition}

\begin{definition}[Hidden state]\label{def:pl:the-research-environment:hidden}
A notebook has \emph{hidden state}\index{hidden state} when its saved outputs depend on something its cells, run from the top, would not
reproduce: a cell run out of order, a cell run twice, a cell edited or deleted after it ran, a name defined only at the prompt.
\end{definition}

The notebook is the fastest way to explore data and the easiest way to publish a number nobody can recompute. Pimentel, Murta, Braganholo and
Freire collected Jupyter notebooks from GitHub and tried to run them: out of 863\,878 attempted executions of valid notebooks, only 24.11\%
executed without errors and only 4.03\% produced the same results. Execution counters give \omterm{def:pl:the-research-environment:hidden}{hidden state} away --- a counter skipped means a
cell re-run or run and deleted, a counter lower than the cell above it means execution out of document order --- but only if something reads
them.

The chapter's planted notebook computes the headline of the hook's debate on a synthetic panel of twenty countries from 1946 to 2009: the
average real growth of country-years with debt above 90\% of output, averaged by country. Its history is written down in the code: the
cells, in document order, and what was run, in the order it was run.

\omcode{code/platforms/15-the-research-environment/python/pl_envaudit.py}{44}{53}{The planted notebook: five code cells in document order,
a cell that was run and then deleted, and the order in which the analyst ran them --- the cleaning cell, near the top, last.}\label{lst:pl:the-research-environment:planted}

\begin{omfigure}
\begin{tikzpicture}[font=\scriptsize,cell/.style={draw,rounded corners=2pt,minimum width=5.6cm,minimum height=0.55cm,anchor=west,
  fill=omMeth!10,font=\ttfamily\scriptsize},>=Stealth]
\foreach \y/\n/\s in {0/1/{df = pd.read\_csv(\textquotedbl{}panel.csv\textquotedbl{})},-0.8/6/{df = df[df.year >= 1950]},-1.6/2/{high = df[df.debt > 90]},
  -2.4/4/{sel = high[\textasciitilde high.country.isin(exclude)]},-3.2/5/{result = sel.groupby(\dots)\dots{} -> 1.96}} {
  \node[cell] at (0,\y) {\s};
  \node[anchor=east,font=\scriptsize] at (-0.1,\y) {[\n]};}
\node[cell,fill=omThm!12,dashed] (del) at (7.2,-2.0) {exclude = [\textquotedbl{}C01\textquotedbl{}, \dots, \textquotedbl{}C05\textquotedbl{}]};
\node[anchor=west,font=\scriptsize\itshape,omThm] at (7.2,-2.6) {run as [3], then deleted};
\node[font=\scriptsize\itshape] at (2.8,0.65) {document order, with execution counters};
\draw[->,omThm,thick] (5.7,-0.8) .. controls (6.6,-1.5) and (6.6,-3.0) .. (5.7,-3.2);
\node[omThm,font=\scriptsize,anchor=west] at (6.5,-3.4) {ran after the result};
\end{tikzpicture}

\omcaption{The planted notebook as saved: code cells in document order with their execution counters. The cleaning cell, second on the
page, ran sixth, after the result; the cell that set \texttt{exclude} ran third and was deleted. The saved 1.96 comes from that
history.}\label{fig:pl:the-research-environment:notebook}
\end{omfigure}

Saved, the notebook shows 1.96 under its result cell. Nothing in the document says that five countries were excluded, or that the years before
1950 were still in the data when the number was computed.

\section{Auditing the document}

\texttt{firm.envaudit} reads the notebook's JSON and nothing else (\cref{lst:pl:the-research-environment:audit}). It walks the code cells
in document order with their counters, and parses each cell to see which names it reads and which it binds.

\omcode{code/firm/envaudit/firm_envaudit.py}{89}{118}{The audit: counters out of document order, skipped counters, outputs without a run,
cells never run, and names read before any cell above binds them --- or bound by no cell at all.}\label{lst:pl:the-research-environment:audit}

On the planted notebook it finds three things (\cref{tab:pl:the-research-environment:audit}). The cell that selects country-years above 90\%
ran second, but a cell above it ran sixth: the cleaning cell was run after the result, so the displayed data are not the data the result
used. Counter~3 is missing: a cell was run and is gone. And the selection cell reads \texttt{exclude}, which no cell in the notebook defines.

\begin{table}[ht]
\centering\small
\begin{tabular}{llp{7.4cm}}
\toprule
finding & cell & detail \\
\midrule
out of order & 3 & counter 2 below the cell above it (6, cell 2) \\
skipped counter & --- & counter 3 missing: a re-run or a deleted cell \\
never defined & 4 & \texttt{exclude} \\
\bottomrule
\end{tabular}
\caption{The audit of the planted notebook, from its JSON alone (cell 0 is the title).}\label{tab:pl:the-research-environment:audit}
\end{table}

Rerun from the top in a fresh interpreter (\cref{lst:pl:the-research-environment:rerun}), the notebook fails: \texttt{NameError: name
'exclude' is not defined}. The number on the page cannot be produced by the page.

\omcode{code/firm/envaudit/firm_envaudit.py}{132}{140}{The rerun: the code cells, top to bottom, as one script, in a fresh interpreter in a
subprocess.}\label{lst:pl:the-research-environment:rerun}

\section{From notebook to library}

\begin{definition}[Research library]\label{def:pl:the-research-environment:library}
A \emph{research library}\index{research library} is the shared, versioned, tested code that research notebooks and pipelines import:
calculations promoted out of notebooks once they are used more than once, with their choices as explicit arguments and their behaviour
pinned by tests.
\end{definition}

The fix is to decide, in the open, what the calculation is. With no exclusion and the cells run as they stand, the notebook prints 2.44. The
calculation belongs in the library (\cref{lst:pl:the-research-environment:library}): the threshold, the first year, the excluded countries and
the weighting become arguments, and each has a test. The library reproduces all three numbers the debate turns on: 1.96 with the five countries
excluded and all years (the saved number, explained), 2.44 with every country from 1950 (the corrected number), and 2.90 if every
country-year counts once instead of every country --- the weighting choice the hook's critique also questioned. The five excluded countries
hold 204 of the 455 country-years above 90\% since 1950.

\begin{omfigure}
\begin{tikzpicture}
\begin{axis}[width=12cm,height=5.6cm,xbar,bar width=7pt,xlabel={average growth above 90\% debt (\%)},xmin=0,xmax=3.6,
  symbolic y coords={all countries; from 1950,all countries; all years,five excluded; from 1950,five excluded; all years},
  ytick=data,enlarge y limits=0.2,tick label style={font=\small},label style={font=\small},grid=major,grid style={gray!20},
  table/col sep=comma,nodes near coords,nodes near coords style={font=\scriptsize,/pgf/number format/fixed,/pgf/number format/fixed zerofill,/pgf/number format/precision=2},
  legend style={font=\footnotesize,at={(0.5,1.03)},anchor=south,legend columns=2}]
\addplot[fill=omDef!45,draw=omDef] table[y=case,x=by_country]{figdata/platforms/15-the-research-environment/choices.csv};
\addlegendentry{mean of country means}
\addplot[fill=omMeth!45,draw=omMeth] table[y=case,x=by_country_year]{figdata/platforms/15-the-research-environment/choices.csv};
\addlegendentry{mean of country-years}
\end{axis}
\end{tikzpicture}

\omcaption{The headline under the three choices the notebook hid: excluding five countries or not, starting in 1946 or 1950, and averaging
countries or country-years. The saved 1.96 is the top pair's mean of country means (blue); the corrected 2.44, the bottom pair's. Data:
\texttt{fig\_envaudit.py}.}\label{fig:pl:the-research-environment:choices}
\end{omfigure}

\omcode{code/platforms/15-the-research-environment/python/research_lib.py}{8}{24}{The promoted function: the choices that lived in the
notebook's memory are arguments.}\label{lst:pl:the-research-environment:library}

\section{Review}

\begin{definition}[Code review]\label{def:pl:the-research-environment:review}
\emph{Code review}\index{code review} is the examination of a change to shared code by someone other than its author before it is merged:
of what the change does, whether its tests show it, and whether it does what its description says.
\end{definition}

Review is where the research review of Book~7 (chapter~1) meets the codebase. A reviewer cannot review a notebook's \omterm{def:pl:the-research-environment:hidden}{hidden state}, only a
diff: the promotion of a calculation to the library is what makes it reviewable. The chapter's promotion checklist, run by
\texttt{promotion\_checklist}, is the minimum a reviewer asks for: the audit finds no \omterm{def:pl:the-research-environment:hidden}{hidden state}, the notebook reruns from the top, the rerun
reproduces the saved result, the library functions are tested, the environment matches the lock, the data match the snapshot.

\begin{definition}[Monorepo]\label{def:pl:the-research-environment:monorepo}
A \emph{monorepo}\index{monorepo} is a single version-control repository holding the code of many projects, libraries and teams, so that a
change to a library and to every caller can be made, reviewed and tested together.
\end{definition}

Google's account of its own (Potvin and Levenberg, 2016) describes a single repository serving as a common source of truth for tens of
thousands of developers, made workable by systems and workflows built for its scale. This series' own code (\texttt{code/firm/}) is a \omterm{def:pl:the-research-environment:monorepo}{monorepo} in miniature: a
component of one book is used by the next, and a change to it is tested against every chapter that imports it.

\section{Environments and containers}

\begin{definition}[Container image]\label{def:pl:the-research-environment:container}
A \emph{container image}\index{container image} is a packaged file system and configuration --- an operating-system layer, the interpreter,
the libraries, the code --- from which identical isolated environments can be started on any machine that runs containers.
\end{definition}

The Open Container Initiative specifies the format: an image manifest, an optional image index, a set of file-system layers and a
configuration, so that tools for building, transporting and running images interoperate. A container freezes an environment; an environment
lock (Book~7, chapter~29) describes one. The chapter uses the lighter of the two: \texttt{capture\_env} lists the interpreter and every
installed distribution and \texttt{compare\_env} checks them against the repository's \texttt{requirements.txt}; the data are checked
against a manifest of hashes (\texttt{snapshot}). Where the lock is enough, a container adds nothing but weight; where system libraries or
compilers matter (chapter~9's extensions), it is the only honest record.

\begin{dated}{2026-09}{Reproducibility tooling}\label{dat:pl:the-research-environment:tools}
The Jupyter notebook format (nbformat) stores each code cell's \texttt{execution\_count} as an integer or null, and its outputs as a list;
the Open Container Initiative's image specification defines an image as a manifest, an optional index, file-system layers and a
configuration. Both are open specifications; notebook linters and container build tools built on them change often, the formats less.
\end{dated}

\section{Reproducibility as a routine}

The last step reproduces the corrected number from a clean directory: the promoted library, the data file, and a three-line script that
imports the library and prints the headline. The data file is checked against the snapshot's hash (no difference), the environment against
the lock (no difference), and the script prints 2.44 in a fresh interpreter. It is the reproducible result of Book~7 (chapter~29), obtained
not by care but by routine: a result is published only with the command that recomputes it from a clean checkout.

\begin{omfigure}
\begin{tikzpicture}[font=\scriptsize,box/.style={draw,rounded corners=2pt,align=center,minimum height=0.9cm,minimum width=2.2cm,
  fill=omDef!10},>=Stealth]
\node[box] (a) at (0,0) {clean\\ directory};
\node[box] (b) at (2.8,0) {data file\\ vs snapshot hash};
\node[box] (c) at (5.6,0) {environment\\ vs lock};
\node[box] (d) at (8.4,0) {fresh interpreter:\\ library + script};
\node[box,fill=omMeth!15] (e) at (11.2,0) {2.44};
\foreach \x/\y in {a/b,b/c,c/d,d/e} {\draw[->] (\x) -- (\y);}
\node[font=\scriptsize\itshape] at (2.8,-0.8) {no difference};
\node[font=\scriptsize\itshape] at (5.6,-0.8) {no difference};
\end{tikzpicture}

\omcaption{Reproduction as a routine: a clean directory with the promoted library and the data file, the data checked against the snapshot's
hash and the environment against the lock, and the headline recomputed in a fresh interpreter.}\label{fig:pl:the-research-environment:repro}
\end{omfigure}

\begin{example}[Rerun from the top]\label{ex:pl:the-research-environment:rerun}
The planted notebook shows 1.96. The audit finds a cell run out of order, a deleted cell and a name defined nowhere; the rerun from the top
fails on that name. Without the exclusion, and with the cells as they stand, the result is 2.44; the library shows that 1.96 is exactly the
average with five countries excluded and the years before 1950 still in the data, and that pooling country-years instead of countries would
give 2.90. From a clean directory, with the data snapshot and the environment lock checked, the library reproduces 2.44.
\end{example}

\section{Tutorial: audit, rerun, promote, reproduce}\label{tut:pl:the-research-environment:tut}

\begin{tutorial}
\textbf{Goal.} Find a notebook's \omterm{def:pl:the-research-environment:hidden}{hidden state} from its JSON, show that its number cannot be recomputed, fix and promote the calculation, and
reproduce it from a clean directory. \textbf{End state:} \cref{tab:pl:the-research-environment:audit} and the reproduction's three checks.
\begin{tutsteps}
\item \textbf{Plant}: \texttt{pl\_envaudit.planted(dir)} writes the data and the notebook with the outputs of its history.
\item \textbf{Audit}: \texttt{firm\_envaudit.audit(nb)}.
\item \textbf{Rerun}: \texttt{firm\_envaudit.rerun(nb, dir)}, then with \texttt{fixed(nb)}.
\item \textbf{Promote}: \texttt{research\_lib.average\_growth} and its tests.
\item \textbf{Reproduce}: \texttt{pl\_envaudit.reproduce(clean\_dir, dir)}: snapshot, lock, rerun.
\end{tutsteps}
\textbf{What to change next.} Run the audit on your own notebooks; add a check that every saved output of a notebook is reproduced by its rerun
(outputs compared cell by cell).
\end{tutorial}

\section{Build: the environment audit}\label{bld:pl:the-research-environment:envaudit}

\begin{build}
\bfield{Purpose} The platform's gate between exploration and anything others rely on: a notebook's number becomes a result only when it can
be recomputed from the top, from a clean checkout, in the locked environment, on the snapshot of its data.
\bfield{Interface} \texttt{load}, \texttt{code\_cells}, \texttt{audit -> [Finding]}, \texttt{saved\_output}, \texttt{rerun -> Rerun},
\texttt{capture\_env}, \texttt{compare\_env}, \texttt{snapshot}, \texttt{check\_snapshot}, \texttt{promotion\_checklist}.
\bfield{Rules} The audit reads the document only; the rerun uses a fresh interpreter in a subprocess, never the author's kernel; environments
are compared with the lock file, data with hashes.
\bfield{Acceptance tests} \texttt{code/firm/envaudit/tests/}: a clean notebook has no findings and reruns; a planted one yields each kind of
finding and a failed rerun; imports, functions and built-ins are not findings; lock and snapshot comparisons; the checklist passes and fails
as it should.
\bfield{Stretch} Cell-by-cell comparison of saved and rerun outputs; a container build from the lock; the audit as a pre-commit hook in the
\omterm{def:pl:the-research-environment:monorepo}{monorepo}.
\end{build}

\begin{omsources}
\item T.~Herndon, M.~Ash and R.~Pollin, ``Does high public debt consistently stifle economic growth? A critique of Reinhart and Rogoff'',
\emph{Cambridge Journal of Economics} 38(2), 2014 (PERI working paper 322, 2013).
\item J.~F. Pimentel, L.~Murta, V.~Braganholo and J.~Freire, ``A large-scale study about quality and reproducibility of Jupyter
notebooks'', MSR 2019.
\item R.~Potvin and J.~Levenberg, ``Why Google stores billions of lines of code in a single repository'', \emph{Communications of the ACM}
59(7), 2016.
\item Jupyter nbformat documentation; Open Container Initiative image specification.
\end{omsources}

\section{Exercises}

\begin{exercise}[$\star$]\label{exo:pl:the-research-environment:1}
A notebook's code cells show counters 1, 2, 5, 3, 6 in document order. What does the audit report?
\end{exercise}

\begin{exercise}[$\star$]\label{exo:pl:the-research-environment:2}
Why does the audit report \texttt{exclude} as never defined rather than used before defined?
\end{exercise}

\begin{exercise}[$\star$]\label{exo:pl:the-research-environment:3}
Which two differences between the saved and the corrected calculation account for 1.96 against 2.44?
\end{exercise}

\begin{exercise}[$\star\star$]\label{exo:pl:the-research-environment:4}
Can a notebook with consecutive counters in document order still have \omterm{def:pl:the-research-environment:hidden}{hidden state}? Give an example.
\end{exercise}

\begin{exercise}[$\star\star$]\label{exo:pl:the-research-environment:5}
Why is weighting by country (2.44) or by country-year (2.90) a choice that belongs in the function's arguments, and which would you report?
\end{exercise}

\begin{exercise}[$\star\star$]\label{exo:pl:the-research-environment:6}
When is an environment lock enough, and when is a \omterm{def:pl:the-research-environment:container}{container image} needed?
\end{exercise}

\begin{exercise}[$\star\star\star$]\label{exo:pl:the-research-environment:7}
\emph{Coding.} Add to \texttt{firm.envaudit} a check that compares each cell's saved output with the rerun's output for that cell. What does
it report on the fixed notebook, and why is it harder than comparing the last output?
\end{exercise}

\begin{exercise}[$\star\star\star$]\label{exo:pl:the-research-environment:8}
\emph{Find the flaw.} ``The notebook is in version control and the reviewer approved it, so its result is reproducible.''
\end{exercise}

\section{Problem: Rerun From the Top}

\begin{problem}[{Weekend problem --- a number the page cannot produce}]\label{pb:pl:the-research-environment:1}
The planted notebook of \cref{lst:pl:the-research-environment:planted}, \texttt{firm.envaudit} and the promoted library.

\medskip\noindent\textbf{Part I --- The notebook.}
\begin{enumerate}
\item What does a Jupyter notebook store for each code cell?
\item In what order were the cells run, and in what order are they shown?
\item What number does the notebook show?
\item What is \omterm{def:pl:the-research-environment:hidden}{hidden state}, in three forms?
\item What does the large-scale study of notebooks on GitHub report?
\end{enumerate}

\medskip\noindent\textbf{Part II --- The audit.}
\begin{enumerate}[resume]
\item What three findings does the audit make, and from what?
\item What does the rerun from the top do?
\item What does each finding explain about the saved number?
\item What does the audit not see?
\item Why rerun in a fresh interpreter in a subprocess?
\end{enumerate}

\medskip\noindent\textbf{Part III --- The library.}
\begin{enumerate}[resume]
\item What becomes an argument of the promoted function?
\item What are the three numbers it reproduces, and under which choices?
\item How many country-years above 90\% since 1950 do the five excluded countries hold?
\item What does a reviewer check before the promotion is merged?
\item Why does a \omterm{def:pl:the-research-environment:monorepo}{monorepo} help?
\end{enumerate}

\medskip\noindent\textbf{Part IV --- The verdict.}
\begin{enumerate}[resume]
\item State the \emph{named result}: the headline number as saved and as rerun, the findings that explain the difference, and the reproduction
of the corrected result from a clean environment.
\item What would you require before a notebook's number goes into a research report?
\item When would you build a \omterm{def:pl:the-research-environment:container}{container image}?
\item What should a research log (Book~7, chapter~1) record about this analysis?
\item In one sentence: what makes a number reproducible?
\end{enumerate}
\end{problem}

\section{Interview questions}

\begin{interviewq}[$\star$ \iqroles{researcher}]\label{iq:pl:the-research-environment:1}
A colleague's notebook shows a result you cannot reproduce. What do you check?
\end{interviewq}

\begin{interviewq}[$\star\star$ \iqroles{researcher, developer}]\label{iq:pl:the-research-environment:2}
When should code move from a notebook to a library?
\end{interviewq}

\begin{interviewq}[$\star\star$ \iqroles{developer}]\label{iq:pl:the-research-environment:3}
What does a container give you that a requirements file does not, and what does it cost?
\end{interviewq}

\begin{interviewq}[$\star\star$ \iqroles{researcher}]\label{iq:pl:the-research-environment:4}
How would you make every number in a research report reproducible a year later?
\end{interviewq}

\begin{interviewq}[$\star\star\star$ \iqroles{developer, researcher}]\label{iq:pl:the-research-environment:5}
Design the research environment for a team of twenty researchers sharing data and code.
\end{interviewq}
